\chapter{Anekdoter}
Vi har fået en person fra hver dekade af \TKETs eksistens til at skrive en lille
tekst om tiden i \TKET . Vi starter med en historie fra en af \TKETs stiftere.
\section{1956-1966}

Jeg havde i 55–56 været stud.act. i København og været medlem af Parentesen, og da jeg syntes, vi burde have en tilsvarende forening i Århus, satte jeg en dag et opslag op om et stiftende møde. De fleste af de ca. 20 studerende og en del af de endnu færre medarbejdere ved fakultetet mødte op, og foreningen blev startet. Den havde ikke noget navn i starten, men i løbet af vinteren udskrev vi en konkurrence om det bedste og bedst begrundede navn - konkurrencen blev vundet af Olaf Pedersen, som kunne opremse en række egenskaber ved et tågekammer, som i overført betydning var positive for foreningen.

Så vidt jeg husker, var jeg formand fra starten i 56 til 59. De første år havde vi en af fakultetets medarbejdere med i bestyrelsen. Axel Voigt var med i den første, og vi holdt bestyrelsesmøder på hans kontor. Da vi forberedte den første fest, snakkede vi om at lave en punch til den, og at det ville blive billigere, hvis alkoholprocenterne til den kom fra ren alkohol. Voigt sagde, at han naturligvis ikke kunne give os noget af sin alkoholforsyning, men at han lige havde et ærinde - så en god portion af den alkohol, Voigt havde stående, kom til at danne grundlag for punchen. Der var i øvrigt ikke nogen af os, der havde erfaring i brygning af punch, så den første, vi lavede, indeholdt en del alkohol, men smagte ikke af noget. Resultatet var, at mange af deltagerne (uerfarne, som vi dengang var) fik ikke bare lidt for meget at drikke. Jeg husker, at en af deltagerne blev kørt hjem i en trillebør af nogle kammerater og læsset ind af vinduet til en kulkælder, hvor han vågnede næste eftermiddag.

Efter den første fest blev jeg inviteret til middag en aften hos en af medarbejderne på Kemisk, Arne Berg. Først hen på aftenen fandt jeg ud af, hvad formålet med invitationen var: Hans kone, der var læge, ville lære mig at lave punch med mindre alkohol og mere smag. For at servere vores form for punch til en fest "det var sådan noget, man fik børn af!". De, der risikerede at blive gravide, må have været vores gæster, for der var (næsten?) kun mænd i den første årgang.

I sæsonen 57–58 lavede vi en revy til en af festerne. Jeg havde skrevet en del af teksten, var instruktør og var en gennemgående person i revyen. Vi havde til denne fest inviteret nogle bestyrelsesmedlemmer fra Parentesen, og udover at lave (og tilsmage) punchen skulle jeg tage mig af dem om eftermiddagen (det indebar bl.a. besøg på Thorups Bodega (eller hed det Thorups vinkælder?) m.m.). Så resultatet var, at jeg ikke kunne huske meget af min tekst om aftenen. Det var ellers en udmærket revy

Selv om jeg mest i det foregående har skrevet om fester, var foreningen i de første år en faglig foredragsforening, og vi startede også festerne med et (festligt) fagligt foredrag. Foredragsholderne var altid lokale (eller folk, der alligevel var i byen – det kostede i hvert tilfælde aldrig noget) – medarbejdere ved institutterne eller lokale gymnasielærere.

\epigraph{Per Mark, FORM 1956–1958}



\section{1966-1976}

Som nyvalgt bestyrelsesmedlem, og guldsmedesøn, var det mig magtpåliggende, at \CERM{}s kæde blev forsynet med en behørigt graveret plade. Denne tradition er vist fortsat (på nye plader). Jeg kan især huske tre arrangementer:

Vi prøvede med en bridge-hyggeaften. Det var dog ikke nogen stor succes, selvom den senere dekan Karl P. deltog. Men øl og bridge er vist ikke en god blanding – især ikke når nogen af deltagerne kommer i forventning om en gang seriøs bridge.

Så gik det straks meget bedre med en Ludo-turnering, hvor brikkerne i finalerunden var øl. Vi havde slet ikke drømt om at det blev en så overvældende succes og at det ligefrem blev en klassiker.

Hypnose var dengang et hot emne og det skulle vi også prøve til en hyggeaften. Vi havde regnet med mange, men at Ali Hamann kunne fylde Aud. E godt og vel, var langt over forventning.

Årets julerevy tog afsæt i en voldsom snestorm i november 1971, som stoppede trafikken over det meste af landet. Folk måtte i stor stil indkvarteres privat, der hvor de sad fast i sneen. De flinke værter lod selvfølgelig folk benytte telefonen, og der blev efterladt mange ubetalte telefonregninger (det var længe før mobiltelefonen). Det gav godt stof til julerevyen. I den efterfølgende påske blæste det igen, så vejene var overstrøede med væltede campingvogne. Det gav endnu et stormfuldt indslag til maj-festen.

\epigraph{Mogens Kamstrup Erlandsen, SEKR 1971–1972}

\section{1976–1986, Fra Crosby til Randers}

Holdet bag julerevyen 1981 var rebelske og var derfor ikke helt tilfredse med, at både White Christmas og cancan’en var obligatoriske. Holdet var selvfølgelig alt for udannede til at kende forskellen på tvang og tradition. Da cancanen jo i sig selv var en stor kunstnerisk udfordring, rettedes oprøret mod sangen. White Christmas udfylder en vigtig funktion, idet hele revyholdet – minus tre – synger fællessang med publikum. Sangteksten er traditionen tro trykt i det udleverede sanghæfte, og den i hast forfattede danske tekst, prisende en årstidsøl fra Thor, blev derfor præsenteret for publikum på en flip-over på scenen. Jubelen var stor.

\subsection{Bundgaard, er du for gammel nu?}

foråret 1982 fik \TKET en henvendelse fra en dengang 50-årig lektor på Matematisk Institut, der på et kraftigste opfordrede \TKET til at spille en rolle i den festdag, der skulle holdes i anledning af Svend Bundgaards 70 års fødselsdag. Dagens arrangement bestod af en række festforelæsninger, men opfordringen lød helt specifikt på at \TKET skulle afbryde selve Svend Bundgaards forelæsning. Formodningen var, at et sådant indslag ville være velkomment i en lang dag, og at det også ville varme fødselarens hjerte at se unge aktive studerende i fuld aktion.

Et hurtigtarbejdende udvalg udvalgte nogle gamle \TK{}-travere, men en egentlig nyforfatning blev også iværksat. Der gælder om mange af revynumrene, at de er skrevet til bestemte aktører, altså rent Pulp Fiction, hvor man skriver en sketch til dem, der skal spille. Re-casting er derfor en svær proces. Måske endda umulig. Samtidig gjaldt det forhold, at flere af \TKETs poeter havde en endog ganske hvas pen. Oden til Bundgaard starter derfor „Jamen Bundgaard, er du for gammel nu?“.

Revyindslaget, der foregik foran tavlerne i Aud. E, varede tæt op mod den halve time. Det samt en ikke helt optimal træning (og timing) og en ganske graverende mangel på fortæppe gjorde, at indsatsen langt fra indfriede forventningerne. Lektorens navn skal derfor forblive usagt. \TKET lader sig aldrig ryste, så oden til Bundgaard fik selvfølgelig repremiere i majrevyen samme år.

\subsection{Kantinens 20 års-jubilæum}

I 1984 afholdt Matematisk Kantine et brag af en jubilæumsfest. Festen var så stor, at hele kantinen var dækket op, og revyen skulle derfor foregå i det fatale Aud. E. Belært af den tidligere indsats i samme lokalitet blev overgangene planlagt med meget stor omhu. Virkemidlet var filmlærredet. Ikke selve lærredet, men den mulighed at placere et par aktører på den smalle afsats ved lærredets bund. Det var så tilpas højt oppe, at når publikum kiggede derop, var der frit lejde på hovedscenen på gulvet i auditoriet. Revyen var selvfølgelig igen århundredets revy, men under den senere indtagelse af handelsrusmidler indvendte Jørgen Andersen, at det tenderede en alvorlig forsømmelse, at hans noget aldrende vise om at bryde det sjette bud ikke havde fundet vej til jubilæumsrepertoiret. Hans indvendinger må have gjort et vist indtryk, for den blev i alt fald fremført i den næstkommende julerevy.

\epigraph{Peter Fristrup, \CERM 1981–1982}

\section{1986–1996, Night on Earth}

\textbf{11h35, TK, Aarhus Universitet}Røde Martin (G\VC) sidder over en cola, FORM har kontortid og er netop luntet i Mat.Kant. for at tilkæmpe sig lidt lasagne, måske med en enkelt frikadelle skjult under sovsen. En kapsel knipses, og svirper forbi G\VC{}’s øre – „Hej \VC “ siger han. \VC griner smørret mens han sætter sig til rette med sin Mat.Kant.-lasagne...

\noindent
\textbf{19h29, TK, Aarhus Universitet} „The Julekalender“ er netop stoppet, folk halser ned fra TV-stuen på matematisk. FORM åbner mødet. Alle får øl. Julefestforberedelserne er på deres højeste.

\noindent
\textbf{19h34, TK, Aarhus Universitet}  G\VC og Ole (Inventar) styrter ind mens de råber på møde-øl. FORM kigger på uret – 19h34 – de får en øl, smiler fjoget.

\noindent
\textbf{20h30, TK, Aarhus Universitet}  Voldsom debat. \VC har knipset en kapsel op i paraplyen, der derefter tumlede ned og ramte Ole (Inventar). Ole siger det gjorde ondt. Ole har været så meget på TK at han tidligere blev udnævnt til Inventar. Inventar må ikke hærges, så \VC er på vej på skyldnertavlen, men argumenterer kraftigt ved at stille en kasse øl på bordet. Debatten forstummer overfor \VC{}’s overbevisende argumenter.

\noindent
\textbf{24h00, Broen over Søen, Aarhus Universitet} En ærværdig gruppe FORMænd klædt i kapper og hætter står med fakler på østsiden. En pjaltet flok fra Næstformationen står på vestsiden. Langsomt nærmer en af de hætteklædte sig broens østside, men uslingene på vestsiden skyder med vandpistoler. En udveksling finder sted. En vigtig task fra en skrivemaskine, mod et smukt rat.

\epigraph{Niels Madsen, FORM 1991–1992, \KASS 1992–1993}

\section{1996–2006, Historien om da \RemToR var til salg}

Ikke så længe efter at vores kære nation vandt EM i en velkendt sportsgren, der har to hold af 11, og noget luft med grisehud omkring som vigtigste elementer, var der en NF, opkaldt efter den lyd man siger når man har glemt at høre sin øl på en torsdag. . . UPS!

Denne ups-NF havde, udfra de kriterier man iøvrigt sammenligner NF’er med, så vidt erindret gjort det ganske glimrende, indtil han fandt det passende at tage et par uger ud at rejse midt i sin embedsperiode?!?

Dette var vi nogle løbeseddelsafhængige kammersugere, der fandt værdigt for en straf. Denne skulle naturligvis bestå i at når NF’en kom hjem var \RemToR væk. . . (vi var vist ikke de første, eller sidste, der lavede det nummer).

Dette var vi nogle løbeseddelsafhængige kammersugere, der fandt værdigt for en straf. Denne skulle naturligvis bestå i at når NF’en kom hjem var \RemToR væk. . . (vi var vist ikke de første, eller sidste, der lavede det nummer).

Men men, i vores Ceres-fremkaldte klarsyn indså vi straks, at inden vi overtalte \RemToR til at gemme sig måtte vi have en plan. Indtagelsen af flydende og hørt inspiration fremkaldte lidt efter en plan; snedig i sin enkelhed. . .\RemToR skulle sættes til salg i Stakbogladen!

Endnu et par kolde Ceres kom på bordkanten, og snart indså vi illegale kombattanter at et passende gemmmested også var en nødvendighed. Hvad var mere


naturligt end at finde en kontorbesiddende person i Næstformationens komplement? Vores første kandidat var med på ideen! Denne person var på daværende tidspunkt desuden med i Fysisk Fredagsbar, og foreslog derfor at deres lagerrum var et godt gemmested.

Der var bare lige det lille problem, at andre medlemmer af samme fredagsbars bestyrelse \textit{ikke} lå i Næstformationens komplement, og kunne opdage maskinen stå der. Men men men, det skulle vise sig slet ikke at være noget problem, da denne tidligere NF slet ikke besad evnen at genkende en \RemToR; selv hvis den stod på en hylde i et lagerrum.

Et par dage efter NF’ens hjemkomst blev Katrin indviet i planen, og \RemToR blev sat ca. midt i bogladen til den nette pris af 1.000.000 kr. med det snedige tvist, at hvis man var NF og kom og spurgte efter den kunne man få den gratis. Dette tvist fremgik dog ikke af prismærkningen, og den fortvivlede NF fandt det passende at tilbyde et par flasker hvidvin i bytte.

NF fik sin \RemToR og Katrin fik noget vin, som hun efterfølgende synes vi skulle have. Det gik vi nødtvunget med til og fandt det passende at bruge en hel torsdag på kammeret på at drikke det. En torsdag hvor NF alias UPS også var til stede og slet ikke forstod hvorfor vi gentagne gange råbte „NF! Skål i hvidvin. . . “

Folk, der dyrker Krystal-massage og stjernetegns-healing, kan sikkert se det skæbnebestemte i, at undertegnede senere blev medlem af den ypperste del af næstformationens komplement, nemlig den højtærede og allestedsnærværende FORMJUNTA. Andre (mig selv inklusiv) ville sikkert bare sige: „Man har jo lov at være heldig.“

At mit og \RemToR{}s spor skulle krydses igen (og igen og igen), ja det er en helt anden historie.

\epigraph{Peter Juhl Christansen, \VC 1996–1997, FORM 1997–1998}

\section{2006-2016}

Et af de spørgsmål, man oftest får som hyggekammerhistoriker, er også et af dem, det oftest er sværest at finde svar på, nemlig: “Hvordan I alverden fandt I på at starte DÈN tradition!?!” Som jeg nu vil illustrere, er vores elskede \TKET nemlig hjemsted for utroligt mange arrangementer, der enten er opstået ved en fejl eller blevet traditioner ved et tilfælde.

I 2008 arrangerede vi et foredrag (om Takt og Tone), som viste sig at være kortere, end hvad vi plejede - så vi besluttede at fylde tiden ud med en omgang Lanciers i solgården. Det blev kun bedre af, at en af vores trofaste hængere havde en mor, der kunne levere undervisning og levende musik. Det var tænkt som en nødløsning for at lappe på et for kort foredrag, mennu runder vi snart 20 år med årlige Lanciers-aftener på kammeret, og egentlig kan man vel kun undre sig over, at det tog os så længe at starte den hyggelige tradition.

Det har nu i mange år været kutyme, at man bunder tre drikkevarer som fejring, når man har afleveret speciale. Da jeg opdagede det, må jeg tilstå, at jeg udbrød noget i stil med: “Hvad?! Hvem i ALVERDEN har været så TÅBELIG at finde på…. Åh nej!!!” For det var cirka på det tidspunkt i sætningen, jeg indså, at ansvaret lå hos… mig selv! Jeg havde fundet på at skrive en specialesang til brug ved relevante torsdagsmøder, og den valgte melodi (Katinka Katinka) rimer temmelig mange steder. For at få det til at passe, endte jeg med linjen: “Nu skal der bundes af det, der’r kreer’t - båd’ stående, liggende og siddende!” I min naive tilstand tænkte jeg bare, at det var en kreativ måde at beskrive på, at man skulle drikke fra sit speciale (en allerede etableret tradition). Den logiske konsekvens af at kombinere kammeret med de linjer indså jeg først år senere.

Det er et velkendt faktum, at FORM er dum. Det blev illustreret på underholdende vis, da Formanden 2011/2012 kom til at starte arrangementet “Hej-dag”. Det kan være svært at holde styr på, hvor mange felter, man skal bruge, når man skal holde billetsalg i en hel uge, og FORM måtte en dag indse, at han havde tegnet et felt for meget. Uden at lade sig gå på af den slags detaljer skrev han kækt “HEJ” i det sidste felt. (Det her er endnu et eksempel på, at selv erfarne tågekammerater kan undervurdere den logiske konsekvens af at kombinere kammeret med noget…) Pludselig var der mennesker, der meldte sig til “Hej-dags-vagten”. FORM tog det med højt humør, satte også sig selv på vagten og tænkte, at han måtte finde på noget fjollet, når dagen nærmede sig. Han endte med at beslutte, at de to vagter var hvert sit hold, der spillede mod hinanden for at få det andet hold til at sige en bestemt sætning flest gange i løbet af dagen. Han udtænkte en mystisk sætning til modstanderne, som det var svært at flette ind i samtaler, og så troede han, den var hjemme. Men han havde undervurderet sin \KASS{}, som spillede på modstanderholdet: FORM dukkede op til Hej-dag og fik sætningen “Skulle vi ikke tage og drikke en GD?” - sammen med en hel flaske Gammel Dansk. Og så måtte han jo vinde, hvis han var villig til det! FORMs stærke konkurrencegen, kombineret med at han havde startet dagen med ondskabsfulde tømmermænd, gjorde ikke oplevelsen mindre sjov for resten af kammeret, og i flere år herefter var det en tradition at holde Hej-dag, hvor man meldte sig til et ukendt arrangement, som FORM i sidste sekund fandt på indhold til.

I starten af 2007 fik nogle hængere en idé til at lave sjov med \CERM{}. De ville bilde hende ind, at hun var ved at glemme en elsket tradition på kammeret, som hun naturligvis havde ansvaret for at videreføre: fejringen af 47. dag - med sanghæfte, særligt kreeret happeningssang og naturligvis indtagning af Carlsberg 47 i rigelige mængder. De udnyttede i den sammenhæng, at den pågældende \CERM har gået direkte fra FU til BEST, så de satsede på, hun ikke kunne huske, om der havde været holdt 47. dag året inden. Så skrev de fluks de sidste mange års happeningssange på vegne af de gamle \CERM{}er

\epigraph{Ninni Duch Schaldemose, SEKR 2007-2008}

